\documentclass[11pt,a4paper]{article}
\usepackage[utf8]{inputenc}
\usepackage[T1]{fontenc}
\usepackage[margin=1in]{geometry}
\usepackage{booktabs}
\usepackage{tabularx}
\usepackage{xcolor}
\usepackage{hyperref}
\hypersetup{colorlinks=true, linkcolor=blue!60!black, urlcolor=blue!60!black}

\title{Comparison of the Molecular Design VAE Pipeline\\
       with Existing De Novo Drug Discovery Tools}
\author{Simarpreet Kaur\\\small SR No: 26727 \and
        M.\ Bhagya Sri\\\small SR No: 27003}
\date{}

\newcommand{\yes}{\checkmark}
\newcommand{\no}{--}
\newcommand{\partsupp}{$\circ$}

\begin{document}
\maketitle

\section*{Abstract}
This document compares the Molecular Design VAE pipeline with publicly
available tools in the de novo drug discovery and protein--ligand modelling
space. The comparison is organised by feature, hardware requirements, and
intended purpose.

\section{Feature Comparison}

\begin{table}[h]
\centering
\small
\setlength{\tabcolsep}{4pt}
\renewcommand{\arraystretch}{1.2}
\begin{tabular}{p{6.5cm} c c c c c c c}
\toprule
\textbf{Feature} & \textbf{Ours} & \textbf{MolMIM} & \textbf{DiffDock} & \textbf{GNINA} & \textbf{AF3} & \textbf{Boltz-2} & \textbf{EquiBind}\\
\midrule
De novo molecule generation         & \yes & \yes & \no & \no & \no & \no & \no \\
Runs locally on consumer hardware   & \yes & \no  & \no & \no & \no & \no & \no \\
CPU-only execution supported        & \yes & \no  & \no & \no & \no & \no & \no \\
Open-source code                    & \yes & \yes & \yes & \yes & \no & \yes & \yes \\
Free, no account required           & \yes & \no  & \yes & \yes & \no & \yes & \yes \\
Molecular docking scoring           & \yes & \partsupp & \yes & \yes & \partsupp & \partsupp & \yes \\
NSGA-II Pareto multi-objective      & \yes & \no  & \no & \no & \no & \no & \no \\
Browser UI with 3D viewer           & \yes & \no  & \no & \no & \no & \no & \no \\
SAR / activity cliffs / MMP         & \yes & \no  & \no & \no & \no & \no & \no \\
Scaffold hopping                    & \yes & \no  & \no & \no & \no & \no & \no \\
Retrosynthesis route                & \yes & \no  & \no & \no & \no & \no & \no \\
PAINS / ADMET filtering             & \yes & \no  & \no & \partsupp & \no & \no & \no \\
SDF / PDB export                    & \yes & \no  & \partsupp & \yes & \yes & \yes & \partsupp \\
\bottomrule
\end{tabular}
\caption{Feature comparison. \yes\ supported, \partsupp\ partial, \no\ not supported.
AF3 = AlphaFold~3.}
\end{table}

\section{Hardware Requirements}

\begin{table}[h]
\centering
\small
\renewcommand{\arraystretch}{1.3}
\begin{tabularx}{\textwidth}{p{3cm} p{4.5cm} c X}
\toprule
\textbf{Tool} & \textbf{Minimum GPU} & \textbf{Student laptop?} & \textbf{Approximate cost}\\
\midrule
This project    & None (CPU sufficient)            & Yes        & Free \\
NVIDIA MolMIM   & A100 / H100 / L40S (data centre) & No         & \$2--4/hr cloud or \$10k+ hardware \\
DiffDock        & RTX 3080+                        & Slow       & Gaming GPU required \\
GNINA           & Any CUDA GPU                     & Partial    & Gaming GPU required \\
AlphaFold 3     & A100 / TPU                       & No         & Cloud only \\
Boltz-2         & GPU recommended                  & Slow       & Gaming GPU required \\
EquiBind        & GPU recommended                  & Slow       & Gaming GPU required \\
\bottomrule
\end{tabularx}
\caption{Hardware requirements summary.}
\end{table}

\subsection*{Note on NVIDIA MolMIM}

NVIDIA's BioNeMo NIM documentation lists supported GPUs as A100 (40\,GB
or 80\,GB), H100, H200, B200, GH200, and L40S. These are data-centre
accelerators, not consumer or gaming GPUs. The full BioNeMo virtual
screening pipeline additionally requires at least 1.3\,TB of fast NVMe
storage. Running MolMIM also requires an NGC account and API key. The
cost is therefore substantial: an A100 80\,GB runs USD \$10{,}000--15{,}000
used, and cloud rental is \$2--4/hr. This places MolMIM outside practical
reach of student researchers without HPC infrastructure.

\section{Discussion}

The de novo molecular generation space contains two clearly distinguishable
categories of tools. The first \emph{generates} new molecules; among
publicly available tools this currently includes only NVIDIA MolMIM and
the present project. The second category --- DiffDock, GNINA,
AlphaFold~3, Boltz-2, EquiBind --- does not generate molecules. These
are docking, scoring, or structure-prediction tools that operate on
user-supplied molecules.

The closest comparator is therefore NVIDIA MolMIM. Both tools encode
molecules into a continuous latent space and decode novel structures
from sampled latent vectors. The key difference is deployment. MolMIM
is designed for high-performance computing environments using A100 or
H100 accelerators; the present pipeline targets consumer hardware.
This trade-off is intentional --- MolMIM trains on substantially more
data and produces a richer latent space, while the present project
trades some diversity for accessibility.

Among the docking-only tools, GNINA and DiffDock could complement rather
than replace this pipeline. GNINA's CNN scoring function performs well
on novel chemotypes --- exactly the molecules this project generates.
DiffDock's blind docking is also well-suited as an optional final step
for the highest-scoring generated candidates.

The unique contribution of this project is the combination of generative
modelling with a complete analysis platform: chemical-space visualisation,
structure--activity relationship analysis, activity-cliff detection,
matched molecular pair analysis, scaffold hopping, retrosynthesis, and
3D structure viewing in a single locally-runnable interface.

\section*{Sources}
\begin{itemize}
  \item DiffDock: Corso, St\"ark, Jing, Barzilay, and Jaakkola (2022)
  \item GNINA: McNutt et al.\ (2021)
  \item AlphaFold~3: Abramson et al.\ (2024)
  \item Boltz-2: MIT (2024)
  \item EquiBind: St\"ark et al.\ (2022)
  \item NVIDIA BioNeMo Framework documentation, accessed 2026
\end{itemize}

\end{document}
